\chapter{总结与不足}\label{chap:conclusion}

\section{实验总结}

本文围绕手写数字识别与垃圾分类两个任务完成了完整的深度学习课程设计。实验首先在 MNIST 上训练 LeNet5，取得 98.60\% 的测试准确率，并保存预训练权重；随后将该权重迁移到 26 类垃圾分类任务，完成从零训练、冻结微调、全量微调、LoRA 微调和强数据增强改进实验。

实验结果表明，垃圾分类从零训练效果很差，仅 3.85\%，说明小规模真实图片数据难以支撑 LeNet5 从随机初始化中直接学习。冻结微调只训练最后分类层，准确率为 15.38\%，说明 MNIST 预训练权重有一定迁移价值，但冻结主干限制了新任务适配能力。LoRA 微调在 \texttt{fc1} 和 \texttt{fc2} 上加入低秩分支，可训练参数为 5106，准确率达到 25.38\%，高于冻结微调。全量微调开放全部参数，准确率达到 64.23\%，是本实验中表现最好的策略。

强数据增强改进实验未能提升性能，准确率从 64.23\% 降至 38.85\%。通过预处理可视化可以看到，随机旋转在 $32\times32$ 灰度图上引入了明显黑色边框，可能形成无关伪特征。因此该实验的主要结论不是“数据增强无效”，而是当前增强算子配置存在改进空间，后续应检查填充值、插值方式和增强顺序。

MobileNetV2 进阶对照实验进一步说明，LeNet5 的局限主要来自输入适配和模型容量，而不是训练流程本身。使用 MindCV 官方预训练 MobileNetV2-1.0、RGB 输入和 $224\times224$ 分辨率后，冻结微调达到 87.31\%，LoRA 达到 89.62\%，全量微调配合 warmup + cosine 学习率衰减达到 91.15\%。这为垃圾分类任务提供了更接近真实应用场景的性能参照。

\section{LeNet5 在垃圾分类上的局限}

LeNet5 能够清晰展示卷积神经网络和迁移学习的基本流程，但在垃圾分类任务上存在明显局限：

\begin{enumerate}[label=(\arabic*)]
  \item \textbf{颜色信息丢失}：垃圾分类图片由 RGB 转为灰度，颜色线索被丢弃，而颜色常常是区分塑料、金属、纸张、果皮等类别的重要依据；
  \item \textbf{输入尺寸过小}：$32\times32$ 输入会压缩真实物体的纹理和局部细节，不利于细粒度分类；
  \item \textbf{模型容量有限}：LeNet5 层数浅、参数少，对 26 类真实场景图片的表达能力不足；
  \item \textbf{源任务与目标任务差异大}：MNIST 是简单手写数字，垃圾分类是真实物体图片，二者分布差异大，因此冻结主干迁移效果有限；
  \item \textbf{LoRA 边际收益有限}：LeNet5 本身参数很少，使用 LoRA 不会像大模型微调那样显著节省显存或训练时间，其主要价值是教学性和实现验证。
\end{enumerate}

\section{后续改进方向}

后续可以从以下方向改进：

\begin{enumerate}[label=(\arabic*)]
  \item 使用三通道颜色输入，设计三通道 LeNet 变体，减少颜色信息损失；
  \item 适当增大输入分辨率，例如 $64\times64$ 或 $96\times96$，保存更多纹理细节；
  \item 在 LeNet5 中加入 Dropout、BatchNorm 或更合理的数据增强策略，控制过拟合；
  \item 修正强数据增强的旋转填充值和插值配置，避免在低分辨率灰度图上产生黑色边框伪特征；
  \item 绘制混淆矩阵，分析易混类别并设计针对性增强；
  \item 继续测试 MobileNetV2、ResNet 等更强模型，与 LeNet5 实验形成多尺度对照。
\end{enumerate}

总体而言，本实验完成了从 MNIST 预训练到垃圾分类微调的完整 LeNet5 流程，并通过冻结、全量、LoRA 和数据增强实验展示了不同训练策略在准确率、可训练参数量、训练时间和显存占用上的差异。
